\documentclass[11pt]{article}

\usepackage{amsmath}
\usepackage{amssymb}
\usepackage{mathtools}
\usepackage[margin=1in]{geometry}
\usepackage{booktabs}
\usepackage{longtable}
\usepackage{array}
\usepackage{xcolor}
\usepackage{enumitem}
\usepackage{hyperref}
\hypersetup{colorlinks=true,linkcolor=blue!50!black,citecolor=blue!50!black,urlcolor=blue!50!black}

\DeclareMathOperator{\tr}{tr}
\DeclareMathOperator{\diag}{diag}
\DeclareMathOperator{\clip}{clip}
\newcommand{\R}{\mathbb{R}}
\newcommand{\Wset}{\mathbb{W}}
\newcommand{\Xset}{\mathbb{X}}
\newcommand{\Uset}{\mathbb{U}}
\newcommand{\norm}[1]{\left\lVert #1 \right\rVert}
\newcommand{\code}[1]{\texttt{#1}}
\newcommand{\yes}{{\color{green!55!black}\textbf{USED}}}
\newcommand{\no}{{\color{red!75!black}\textbf{NOT USED}}}
\newcommand{\concept}{{\color{orange!85!black}\textbf{CONCEPT ONLY}}}

\title{\textbf{Mathematical Linkage Audit:\\[2pt]
What of Phase 1 (Mathematical Foundations) Is Actually Used in Phase 3 (System Integration)}}
\author{Companion analysis for the \code{ugv\_uav\_control} package}
\date{\today}

\begin{document}
\maketitle

\begin{abstract}
This document goes through the \emph{complete} mathematics of Phase~1
(\code{01\_Mathematical\_Foundations}) and determines, object by object, how much of it is
carried into Phase~3 (\code{03\_System\_Integration\_ROS2}). The conclusion is sharp and important
for paper writing: the link between the two phases is \textbf{architectural and conceptual, not
numerical}. Phase~3 inherits the \emph{form} of Tube~MPC from Phase~1 --- a nominal plan plus an
ancillary error-feedback law $u=\bar u-Ke$, an invariant-set membership test used as a replan
trigger, and the surrounding RPI / constraint-tightening vocabulary. But \emph{none} of Phase~1's
computed quantities --- the LQR gain $K$, the Riccati matrix $P$, the minimal robust positively
invariant (mRPI) set $\Omega$, the tightened constraints $\bar{\Xset},\bar{\Uset}$, the disturbance
set $\Wset$, or even the plant matrices $A_d,B_d$ --- reach Phase~3. The root cause is a change of
plant model (a $6$-state double integrator in Phase~1 versus a $3$-state single integrator in
Phase~3), which makes Phase~1's artifacts dimensionally incompatible with Phase~3, compounded by
hard-coded gains in the Phase~3 code. Consequently the README claim that the Full System
``leverages the RPI Set and Constraint Tightening derived in Phase~1'' is \textbf{not supported by
the code} and should be reworded.
\end{abstract}

\tableofcontents

\section{Scope and method}
Phase~1 is, mathematically, a single self-contained object: a robust Tube~MPC design for the UAV
modelled as a double integrator. It computes an LQR feedback gain, a disturbance set, an mRPI set,
and tightened constraints offline (\code{simulations/tube\_calculation.py}), saves them to
\code{tube\_data.npz}, and consumes them online in three solver classes
(\code{lcs/mpc\_solvers.py}: \code{StandardMPC}, \code{TubeMPC}, \code{ETTubeMPC}).
It contains \emph{no} EKF, \emph{no} unicycle model, \emph{no} belief-space cost, and \emph{no}
view-quality metric.

Phase~3 is the integrated ROS\,2 mission. Its UAV controller (\code{uav\_mpc\_node.py} planner +
\code{uav\_ancillary\_node.py} tracker) is the only part that overlaps Phase~1's domain. The
EKF (\code{ekf\_node.py}), the belief-space UGV MPC (\code{belief\_mpc\_node.py}) and the
view-quality metric (\code{view\_quality.py}) are Phase-3-native and have no Phase-1 ancestor; they
are out of scope here except where noted.

For each Phase~1 mathematical object below we give the derivation as implemented, then state its
fate in Phase~3 with file/line evidence.

\section{The complete mathematics of Phase 1}

\subsection{Plant: discrete double integrator}
State $x=[p^\top,v^\top]^\top\in\R^6$ with position $p=[p_x,p_y,p_z]^\top$ and velocity
$v=[v_x,v_y,v_z]^\top$; input $u=a=[a_x,a_y,a_z]^\top\in\R^3$ is gravity-compensated acceleration;
sample time $\Delta t=0.05$\,s. The exact discretisation is
\begin{equation}
A_d=\begin{bmatrix}I_3 & \Delta t\,I_3\\ 0 & I_3\end{bmatrix},\qquad
B_d=\begin{bmatrix}\tfrac12\Delta t^2 I_3\\ \Delta t\,I_3\end{bmatrix},\qquad
x_{k+1}=A_dx_k+B_du_k+w_k .
\label{eq:p1plant}
\end{equation}
(\code{tube\_calculation.py:7--23}.)

\subsection{Feedback gain via LQR / DARE}
With $Q=\diag(10,10,10,1,1,1)$ and $R=I_3$, let $P\succ0$ solve the discrete algebraic Riccati
equation
\begin{equation}
A_d^\top P A_d - P - A_d^\top P B_d\,(R+B_d^\top P B_d)^{-1}B_d^\top P A_d + Q = 0,
\end{equation}
and define the gain and closed-loop error matrix
\begin{equation}
K=(R+B_d^\top P B_d)^{-1}B_d^\top P A_d\in\R^{3\times6},\qquad A_K=A_d-B_dK .
\label{eq:p1K}
\end{equation}
$A_K$ is Schur (spectral radius $\rho\approx0.933$). (\code{tube\_calculation.py:31--50}.)
The same $P$ later serves as the MPC terminal cost.

\subsection{Constraints}
Input box (a $30^\circ$ tilt limit, $g\tan30^\circ\approx5.66$, plus a vertical bound) and a state
box matching the lab arena:
\begin{equation}
\Uset:\ |a_x|,|a_y|\le5.66,\ |a_z|\le4.0;\qquad
\Xset:\ p_x\!\in\![-1,1],\,p_y\!\in\![-1.5,1.5],\,p_z\!\in\![0,2.8],\,v\!\in\![-1.5,1.5]^3,
\end{equation}
each converted to halfspace form $H_xx\le h_x$, $H_uu\le h_u$.
(\code{tube\_calculation.py:56--77}.)

\subsection{Disturbance set $\Wset$ (measurement noise through feedback)}
This is the most careful step in Phase~1. Acceleration disturbance $d\in\square([1,1,1.5])$ maps to
state space as $w_{\mathrm{proc}}=B_d d$. Measurement noise $\nu\in\square([0.05,0.05,0.1,0,0,0])$
does \emph{not} enter the plant directly; it corrupts the feedback. With the control law
$u=\bar u-K(x+\nu-\bar x)$ the error $e=x-\bar x$ obeys
\begin{equation}
e_{k+1}=A_Ke_k+\underbrace{w_{\mathrm{proc}}-B_dK\nu_k}_{=:w_k\in\Wset},
\label{eq:p1err}
\end{equation}
so the effective disturbance set is the convex hull
$\Wset=\{\,w_{\mathrm{proc}}-B_dK\nu\,\}$ over the vertices of both boxes
(\code{tube\_calculation.py:93--120}).

\subsection{Minimal robust positively invariant (mRPI) set $\Omega$}
For \eqref{eq:p1err} with $A_K$ Schur and $w_k\in\Wset$, the mRPI set is the convergent Minkowski
series
\begin{equation}
\Omega_\infty=\bigoplus_{i=0}^{\infty}A_K^i\,\Wset .
\end{equation}
Phase~1 computes an outer approximation by the \emph{support-function} method over $72$ directions
(axis-aligned plus pairwise diagonals in $\R^6$), truncated at $N=100$ and scaled by
$1/(1-\rho^{N+1})$ to remain an outer bound, giving a polytope $\Omega_H x\le\Omega_h$.
(\code{tube\_calculation.py:149--234}.)

\subsection{Constraint tightening (Pontryagin difference)}
The nominal trajectory is planned in the shrunken sets
\begin{equation}
\bar{\Xset}=\Xset\ominus\Omega,\qquad \bar{\Uset}=\Uset\ominus K\Omega,
\label{eq:p1tighten}
\end{equation}
computed by subtracting, from each halfspace right-hand side, the support of $\Omega$ (resp.\
$K\Omega$) in that face normal. This guarantees $x_k=\bar p_k+e_k\in\bar{\Xset}\oplus\Omega\subseteq
\Xset$. (\code{tube\_calculation.py:236--272}.)

\subsection{Nominal MPC and ancillary law (as implemented in \code{mpc\_solvers.py})}
\code{TubeMPC} ($N=20$, $\Delta t=0.05$) solves the quadratic program
\begin{equation}
\min_{\bar X,\bar U}\ \sum_{k=0}^{N-1}\norm{\bar x_k-x_{\mathrm{ref}}}_Q^2+\norm{\bar u_k}_R^2
+\norm{\bar x_N-x_{\mathrm{ref}}}_P^2
\end{equation}
subject to the nominal dynamics $\bar x_{k+1}=A_d\bar x_k+B_d\bar u_k$, the \emph{tightened}
constraints \eqref{eq:p1tighten}, and the \textbf{tube initial-condition constraint}
\begin{equation}
\Omega_H\,(x_0-\bar x_0)\le\Omega_h
\label{eq:p1ic}
\end{equation}
(solved with CVXPY/OSQP). The applied input is the ancillary law
\begin{equation}
u=\bar u_0-K(x_0-\bar x_0),\ \text{clipped to }[u_{lb},u_{ub}] .
\label{eq:p1anc}
\end{equation}
(\code{mpc\_solvers.py:146--273}.)

\subsection{Event-triggered variant}
\code{ETTubeMPC} re-solves only when the position error exceeds a threshold for
\code{consecutive\_threshold}$=3$ steps (default trigger $0.15$\,m). (\code{mpc\_solvers.py:316--349}.)

\section{What Phase 3 actually implements (UAV side)}

\subsection{Plant: single integrator}
The Phase~3 UAV is a \textbf{point-mass single integrator} with \emph{velocity} command
$u=[v_x,v_y,v_z]^\top$ and $\Delta t=0.2$\,s (the Gazebo velocity plugin closes the inner loops):
\begin{equation}
p_{k+1}=p_k+\Delta t\,u_k+w_k\quad\Longrightarrow\quad A_d=I_3,\ \ B_d=\Delta t\,I_3 .
\label{eq:p3plant}
\end{equation}
(\code{uav\_mpc\_node.py:41--42,97}.) This is a $3$-state model, not the $6$-state model of
\eqref{eq:p1plant}.

\subsection{Ancillary law and ellipsoidal tube}
\code{uav\_ancillary\_node.py} applies $u=\bar u-Ke$ and an \emph{ellipsoidal} tube test
\begin{equation}
u_k=\bar u_k-K e_k\ (\text{clipped to }[-1,1]^3),\qquad
\text{inside tube}\iff e_k^\top P_0\,e_k\le\alpha_k^2,
\label{eq:p3anc}
\end{equation}
with the EKF-driven adaptive radius
\begin{equation}
\alpha_k=\alpha_{\min}+(\alpha_{\max}-\alpha_{\min})\,
\clip\!\Big(\tfrac{\tr\Sigma_{xy}-\tr_{\min}}{\tr_{\max}-\tr_{\min}},0,1\Big),
\quad \alpha_{\min}=0.5,\alpha_{\max}=2,\tr_{\min}=0.005,\tr_{\max}=0.15 .
\label{eq:p3alpha}
\end{equation}
Crucially the gains are \textbf{hard-coded}:
\begin{equation}
K=\diag(0.6,0.6,0.8),\qquad P_0=0.186\,I_3
\end{equation}
(\code{uav\_ancillary\_node.py:57--58,184--191}). With these, $A_K=I_3-\Delta t\,K=
\diag(0.88,0.88,0.84)$ is Schur --- so the Schur \emph{requirement} of Phase~1 is satisfied, but by
independently chosen numbers.

\subsection{Planner}
\code{uav\_mpc\_node.py} ($N=10$, $\Delta t=0.2$) solves with SLSQP a \emph{nonlinear} cost mixing
goal tracking, an altitude regulariser, a visibility heuristic
$\norm{p_{xy}-g_{xy}}^2/(z^2+0.1)$, and control effort --- with velocity-box bounds and a linear
altitude constraint $z_{\min}\le z\le z_{\max}$, but \textbf{no} terminal cost, \textbf{no}
tightened constraints, and \textbf{no} tube initial-condition constraint
(\code{uav\_mpc\_node.py:92--155}).

\subsection{Event triggers}
Replanning is governed by four prioritised conditions (tube-exit $1a$; ArUco covariance spike
$1b$; curvature anticipation $2$; drift maintenance $3$) with $1$\,s time-lockouts, plus plan
expiry --- not by a consecutive-violation counter (\code{uav\_ancillary\_node.py:198--254}).

\section{Linkage map: object by object}

\renewcommand{\arraystretch}{1.35}
\begin{longtable}{p{4.5cm} p{4.3cm} p{4.3cm} c}
\toprule
\textbf{Phase 1 object} & \textbf{Phase 1 (math)} & \textbf{Phase 3 (code)} & \textbf{Link} \\
\midrule
\endhead
Plant $A_d,B_d$ & Double integrator, $6$ states, accel input, $\Delta t{=}0.05$, Eq.\,\eqref{eq:p1plant} & Single integrator, $3$ states, velocity input, $\Delta t{=}0.2$, Eq.\,\eqref{eq:p3plant} & \no \\
Feedback gain $K$ & LQR/DARE, $3\times6$, Eq.\,\eqref{eq:p1K} & Hard-coded $\diag(0.6,0.6,0.8)$, $3\times3$ & \no \\
Terminal cost $P$ & DARE solution, $6\times6$ & none (planner has no terminal cost) & \no \\
Disturbance set $\Wset$ & $w_{\mathrm{proc}}-B_dK\nu$, Eq.\,\eqref{eq:p1err} & none; surrogate $\alpha_k(\tr\Sigma)$, Eq.\,\eqref{eq:p3alpha} & \no \\
mRPI set $\Omega$ & support-function polytope & ellipsoid $e^\top P_0e\le\alpha^2$, Eq.\,\eqref{eq:p3anc} & \concept \\
Constraint tightening & $\Xset\ominus\Omega,\ \Uset\ominus K\Omega$, Eq.\,\eqref{eq:p1tighten} & none (raw velocity box + $z$ bound) & \no \\
Tube IC constraint & $\Omega_H(x_0-\bar x_0)\le\Omega_h$, Eq.\,\eqref{eq:p1ic} & none & \no \\
Nominal optimiser & QP, CVXPY/OSQP, quadratic & SLSQP, nonlinear heuristic cost & \no \\
Ancillary law $u=\bar u-Ke$ & Eq.\,\eqref{eq:p1anc} & Eq.\,\eqref{eq:p3anc}, same form, diff.\ $K$/dim & \concept \\
Schur stability of $A_K$ & $\rho\approx0.933$ & $A_K=\diag(0.88,0.88,0.84)$ & \concept \\
Event-triggered re-solve & error threshold + counter & ellipsoid test + 4 levels + lockouts & \concept \\
\bottomrule
\end{longtable}

\paragraph{Summary count.} Of the eleven Phase-1 mathematical objects, \textbf{zero} are used
numerically in Phase~3; four cross over as \emph{concepts/structure only} (the ancillary law form,
the invariant-set membership idea, the Schur design requirement, and event-triggering); the
remaining seven --- including everything that constitutes the robustness \emph{guarantee} (the mRPI
set, the disturbance set, and constraint tightening) --- have no Phase-3 implementation.

\section{Why the quantitative link is severed: the model change}
The single decisive fact is the change of plant. Phase~1 derives $K\in\R^{3\times6}$, $P\in\R^{6\times6}$,
$\Omega\subset\R^6$ and $\bar{\Xset},\bar{\Uset}$ for a $6$-dimensional acceleration-controlled double
integrator. Phase~3's UAV state is $3$-dimensional and velocity-controlled. Phase~1's $K$ multiplies
a $6$-vector and returns an acceleration; Phase~3 needs a map from a $3$-vector error to a velocity
correction. The objects are therefore \textbf{dimensionally incompatible} --- Phase~1's outputs
cannot be loaded into Phase~3 even in principle, regardless of code wiring. Everything quantitative
follows from this: with a different $A_K$ the mRPI set differs; without $\Omega$ there is no
Pontryagin tightening; without tightening there is no recursive-feasibility guarantee to inherit.

\section{Code-level findings (relevant to honest reporting)}
\begin{enumerate}[leftmargin=1.4em]
\item \textbf{The tube-parameter load is dead code.} \code{uav\_ancillary\_node.py:47--58} attempts
to load \code{K,P\_0} from a personal absolute path
(\code{/home/shlok-mehndiratta/.gemini/antigravity/scratch/tube\_params.npz}), then on lines
$57$--$58$ --- \emph{outside} the \code{if/else} --- unconditionally overwrites both with
$\diag(0.6,0.6,0.8)$ and $0.186\,I_3$. Even if the file existed, the loaded values are discarded.
The path is also non-portable and is \emph{not} Phase~1's \code{tube\_data.npz}.
\item \textbf{Phase 1's \code{tube\_data.npz} is never referenced by Phase 3.} The only consumers of
\code{tube\_data.npz} are Phase~1's own nodes/sims; a repository-wide search finds no Phase~3 file
loading it.
\item \textbf{Internal inconsistency in the theory companion.} \code{theory\_report.pdf}
(\S6.1) states the fallback gain as $K=2I_3$ giving $A_K=0.6I_3$, and $P_0=6.86\,I_3$. The running
code uses $K=\diag(0.6,0.6,0.8)$ ($A_K=\diag(0.88,0.88,0.84)$) and $P_0=0.186\,I_3$. If the report
is cited in the paper, these numbers should be reconciled with the code.
\item \textbf{Altitude bound nuance.} The theory companion's deviation list says altitude limits are
``declared but never enforced''; in fact \code{uav\_mpc\_node.py:137--145} builds a
\code{LinearConstraint} enforcing $z_{\min}\le z\le z_{\max}$ in the planner. Worth correcting if
that text is reused.
\end{enumerate}

\section{Conclusions and recommendations for the paper}
\begin{enumerate}[leftmargin=1.4em]
\item \textbf{Do not claim Phase~3 ``uses the RPI set / constraint tightening derived in Phase~1.''}
It does not; the robustness machinery of Phase~1 has no Phase-3 implementation, and the gains are
hard-coded.
\item \textbf{Correct framing.} Phase~1 establishes the robust-MPC \emph{formalism} (RPI set,
constraint tightening, ISS of the error dynamics) that \emph{motivates and justifies} the
tube-tracking \emph{architecture} realised in Phase~3 --- i.e.\ theory~$\rightarrow$~architecture,
not theory~$\rightarrow$~numbers. The genuinely shared content is the ancillary law $u=\bar u-Ke$,
the invariant-set membership test as a replan trigger, and the Schur/event-triggering design
principles.
\item \textbf{State the model change explicitly.} The double-integrator~$\to$~single-integrator
switch is defensible (the velocity plugin closes the inner loop) but is the reason Phase~1's
numerics do not transfer; control reviewers will look for this.
\item \textbf{Locate the novelty correctly.} The real Phase-3 contributions are the Phase-3-native
mathematics --- the view-quality-coupled measurement noise $R_k$, the belief-space UGV MPC with
covariance roll-out and FoV gating, and the EKF-covariance-driven adaptive tube radius $\alpha_k$
--- none of which descend from Phase~1.
\end{enumerate}

\end{document}
